\apendice{TODO: Documentación de usuario}

\section{Introducción}

El manual de usuario definitivo se redactará cuando la interfaz esté implementada y validada. Se mantiene este esqueleto para asegurar que la entrega incluya instrucciones coherentes con las pantallas reales.

\section{Requisitos de usuario}

Se indicarán los navegadores soportados, las condiciones de conexión, los límites de descarga y los requisitos para cuentas, correo y funciones administrativas.

\section{Acceso e instalación}

Al tratarse de una aplicación web, el usuario normal no necesitará instalar Batch Downloader. Esta sección explicará cómo acceder al servicio y, si se distribuye una versión local, cómo iniciarla mediante contenedores.

\section{Uso del catálogo}

Se documentarán búsqueda, filtros, selección, tamaño estimado, revisión de fuentes, creación del trabajo, progreso, cancelación y descarga del ZIP.

\section{Bundles y estrellas}

Se explicará la diferencia entre personalizar temporalmente, clonar y editar un bundle propio; la visibilidad pública o privada; el enlace compartible; y la estrella binaria.

\section{Búsqueda semántica}

Se incluirán ejemplos de consultas por intención, explicación de resultados y advertencias sobre el carácter asistido de las recomendaciones.

\section{Administración}

El manual describirá alta y edición de aplicaciones, configuración de fuentes, prueba de resolvers, revisión de solicitudes, logs y estados del catálogo.

\section{Errores frecuentes y soporte}

Se añadirán mensajes reales, causas, pasos de recuperación y forma de comunicar una incidencia mediante su identificador. No se incorporarán capturas ni instrucciones hasta que las pantallas sean definitivas.
